\documentclass[11pt]{article}
\usepackage[T1]{fontenc}
\usepackage{amsmath,amssymb,amsthm,booktabs,geometry,microtype,hyperref}
\geometry{margin=1in}
\hypersetup{hidelinks}
\title{Ellipse Perimeter: Ramanujan Approximation and N-Series Geometric Acceleration\\\large A Chord-Cost and Sinc-Correction Comparison}
\author{C. Wayne Baker}
\date{Original: June 1, 2026\\Revised for review: October 24, 2026}
\begin{document}
\maketitle
\begin{abstract}
We compare Ramanujan's second closed-form \emph{ellipse-perimeter} approximation with inscribed equal-parameter polygons and their scale-extrapolated refinements. The three-scale geometric N-series operator cancels the $N^{-2}$ and $N^{-4}$ perimeter errors and converges at order $N^{-6}$ for each fixed nondegenerate ellipse. We distinguish its base resolution $N$ from its direct chord-evaluation cost $7N$. An exact chord factorization proved in a subsequent companion paper yields the single-polygon sinc-corrected approximation $P_N/\operatorname{sinc}(\pi/N)$, removing the entire algebraic chord-error factor. The remaining periodic midpoint-quadrature error decays exponentially with $N$ for fixed axis ratio $b/a>0$. We reproduce the earlier high-precision comparisons, explain the near-universal normalized sixth-order coefficient, and give chord-count-aware tests of sinc correction. The three mechanisms are mathematically distinct: a closed-form eccentricity approximation, polynomial Richardson cancellation, and analytic chord normalization. No numerical superiority over established elliptic-integral algorithms is claimed.
\end{abstract}
\section{Exact perimeter and Ramanujan's approximation}
For semiaxes $a\ge b>0$, the exact perimeter is
\begin{equation}\label{eq:exact}
P=4aE(e),\qquad e=\sqrt{1-b^2/a^2},\qquad E(e)=\int_0^{\pi/2}\sqrt{1-e^2\sin^2 t}\,dt.
\end{equation}
Ramanujan's second ellipse-perimeter approximation is
\begin{equation}\label{eq:ram}
R_2=\pi(a+b)\left(1+\frac{3h}{10+\sqrt{4-3h}}\right),\qquad h=\left(\frac{a-b}{a+b}\right)^2.
\end{equation}
This is an approximation to the ellipse perimeter, \emph{not} a Ramanujan series for $\pi$. The signed error is positive at sufficiently small $h$: $P-R_2=\pi(a+b)(3h^5/131072+O(h^6))$. In the comparison tables we report \emph{absolute} errors $|P-\widetilde P|$ in the units of $a$ and $b$.
\section{Equal-parameter polygons and scale cancellation}
Set $q(t)=(a\cos t,b\sin t)$ and $t_j=2\pi j/N$. The perimeter $P_N$ is the sum of Euclidean chord lengths $|q(t_{j+1})-q(t_j)|$, with $t_N=2\pi$. For a fixed nondegenerate ellipse the symmetric discretization admits an even-power asymptotic expansion
\begin{equation}\label{eq:exp}
P_N=P+C_2N^{-2}+C_4N^{-4}+C_6N^{-6}+O(N^{-8}).
\end{equation}
Two-scale cancellation gives $(4P_{2N}-P_N)/3=P+O(N^{-4})$. The original three-scale N-series estimator is
\begin{equation}\label{eq:three}
T_N=\frac{P_N}{45}-\frac{4P_{2N}}9+\frac{64P_{4N}}{45}=P+O(N^{-6}).
\end{equation}
Indeed its weights sum to $1$ and its $N^{-2}$ and $N^{-4}$ moments vanish:
\[
\frac1{45}-\frac{4}{9\cdot 4}+\frac{64}{45\cdot16}=0,\qquad
\frac1{45}-\frac4{9\cdot16}+\frac{64}{45\cdot256}=0.
\]
\paragraph{Cost accounting.} The three polygon sums use $N+2N+4N=7N$ chord-length evaluations if computed separately; thus a row headed $N=64$ represents a direct cost of $448$ chords, not $64$. Although sample points can sometimes be reused, the polygon chord lengths for different step sizes are not identical. We use chord counts as a transparent geometric operation proxy, \emph{not} as a measured wall-clock or arithmetic-operation benchmark.
\section{Exact sinc correction and the sixth-order coefficient}
Let $v(t)=\sqrt{a^2\sin^2t+b^2\cos^2t}$ and define the midpoint sum
\[
M_N=\frac{2\pi}{N}\sum_{j=0}^{N-1}v\left(\frac{2\pi(j+1/2)}N\right).
\]
The ellipse chord from $t_j$ to $t_{j+1}$ has length $2\sin(\pi/N)v(2\pi(j+1/2)/N)$, exactly. Hence the factorization from the companion ellipse study \cite{BakerFactorization} is
\begin{equation}\label{eq:sinc}
P_N=\operatorname{sinc}(\pi/N)M_N,\qquad \operatorname{sinc}(x)=\frac{\sin x}{x}.
\end{equation}
Accordingly the single-polygon corrected value is
\begin{equation}\label{eq:correct}
S_N=\frac{P_N}{\operatorname{sinc}(\pi/N)}=M_N.
\end{equation}
The speed $v$ is analytic in each strip $|\operatorname{Im}z|<\rho_*$, where $\rho_*=\operatorname{artanh}(b/a)$ for $0<b<a$. Standard periodic quadrature estimates imply $S_N-P=O(e^{-\rho N})$ for any fixed $0<\rho<\rho_*$; for odd $N$ the $\pi$-periodicity gives the sharper bound $O(e^{-2\rho N})$. For a circle the midpoint rule is exact. These bounds are for a \emph{fixed} ellipse; the constants are not uniform as $b/a\to0$.

Since $M_N-P$ is exponentially small and
\[
\operatorname{sinc}(\pi/N)=1-\frac{\pi^2}{6N^2}+\frac{\pi^4}{120N^4}-\frac{\pi^6}{5040N^6}+O(N^{-8}),
\]
the algebraic sixth-order term of $T_N-P$ is
\begin{equation}\label{eq:six}
T_N-P=-\frac{P\pi^6}{5040\cdot64\,N^6}+O(N^{-8})+O(e^{-\rho N}),
\end{equation}
because $\tfrac1{45}-\tfrac4{9\cdot2^6}+\tfrac{64}{45\cdot4^6}=\tfrac1{64}$. Consequently, $N^6|T_N-P|/P\to\pi^6/(5040\cdot64)$. This explains the nearly uniform small errors across ellipse shapes at sufficiently large $N$: they are proportional to the respective exact perimeters. At modest $N$ and high eccentricity the exponential midpoint contribution can dominate; a universal asymptotic prediction is then premature.
\section{Numerical comparison and precision}
The historical table and initial comparison were checked through AI-assisted high-precision recomputation; the newly added equal-cost $S_{448}$ entries were separately recomputed with 170-decimal-digit working precision using the formulas stated here. These are computational cross-checks, not independent human peer review. The original paper did not specify the precision used to generate its historical table; we therefore do \emph{not} retroactively assert a particular original working precision. Reproduction should use arbitrary precision with at least 60 decimal digits, evaluating \eqref{eq:exact}--\eqref{eq:correct} without binary64 rounding of intermediate chord sums. Errors around $10^{-17}$ in perimeters of order five cannot reliably be resolved by subtracting two ordinary double-precision approximations.
\begin{table}[ht]
\centering\small
\begin{tabular}{@{}rrrrr@{}}\toprule
$b$ ($a=1$) & Base $N$ & Direct chords $7N$ & $|P-R_2|$ & $|P-T_N|$\\\midrule
0.90&256&1792&$2.24\times10^{-17}$&$6.32\times10^{-17}$\\
0.75&256&1792&$4.61\times10^{-13}$&$5.85\times10^{-17}$\\
0.50&256&1792&$2.21\times10^{-9}$&$5.13\times10^{-17}$\\
0.25&256&1792&$1.07\times10^{-6}$&$4.54\times10^{-17}$\\\bottomrule
\end{tabular}
\caption{Absolute perimeter errors for Ramanujan's closed form and the original three-scale method; $N$ is the base resolution, not the chord cost.}\label{tab:old}
\end{table}
For $b=0.50$, the original three-scale errors at base $N=64,128,256$ are $2.10\times10^{-13}$, $3.28\times10^{-15}$, and $5.13\times10^{-17}$, approaching the factor-of-$64$ improvement per doubling predicted by \eqref{eq:six}. At $b=0.25$, base $N=32$ gives $1.81\times10^{-7}$, much larger than the sixth-order algebraic prediction because the quadrature remainder is still material.
\begin{table}[ht]
\centering\small
\begin{tabular}{@{}rrrrrrr@{}}\toprule
$b$ & $|P-R_2|$ & Base $N$ & $7N$ & $|P-T_N|$ & $|P-S_{128}|$ & $|P-S_{448}|$\\\midrule
0.50&$2.2\times10^{-9}$&64&448&$2.1\times10^{-13}$&$1.4\times10^{-33}$&$9.99\times10^{-111}$\\
0.25&$1.1\times10^{-6}$&64&448&$1.8\times10^{-11}$&$2.2\times10^{-17}$&$1.07\times10^{-53}$\\\bottomrule
\end{tabular}
\caption{Comparison including $S_{128}$ (128 chords) and the exactly chord-count-matched $S_{448}$ (448 chords). Rounded absolute errors; the $S_{448}$ values were recomputed at 170-digit precision.}\label{tab:new}
\end{table}
For $b=0.50$, even $S_{64}$ is reported to attain an absolute error about $7.6\times10^{-18}$ from only $64$ chords. Sinc correction is therefore preferable to the three-scale ladder for these particular numerical cases. Near the degenerate limit $b/a\to0$, however, the analyticity strip collapses and $S_N$ need not yet be exponentially accurate at a given finite $N$. Moreover, \eqref{eq:correct} uses an analytic factor containing $\pi$; it is not a claim that straightedge-and-compass operations alone exactly correct the polygon.
\section*{AI-assistance and numerical provenance}
AI tools, including ChatGPT and Claude, assisted with editorial review, formula checks, and numerical comparisons. The earlier tables reflect reported high-precision AI-assisted checks, while the equal-cost $S_{448}$ entries were freshly calculated for this revision with 170-decimal-digit arithmetic. These computational checks do not constitute independent human refereeing. No new empirical runtime measurements are claimed.
\section{Interpretation and conclusion}
The original N-series experiment correctly demonstrates sixth-order Richardson cancellation and shows that sufficient polygonal refinement can exceed the accuracy of Ramanujan's single closed-form ellipse approximation for the tested axis ratios. Explicit cost accounting sharpens that finding: a three-scale estimator at base $N$ uses $7N$ chords, whereas a single sinc-corrected $N$-gon uses $N$ chords. The later exact factorization \eqref{eq:sinc} explains why correcting the chord factor all at once can be far more accurate for an analytic ellipse speed. It also gives a precise interpretation of the original $N=256$ table through the universal algebraic residual \eqref{eq:six}.

The three methods answer different questions. Ramanujan gives a striking fixed algebraic approximation along the eccentricity axis; N-series extrapolation cancels geometric resolution errors without relying on special ellipse factorization; sinc correction exploits exactly such a factorization and leaves periodic quadrature error. These results are not benchmarks against direct elliptic-integral evaluation, AGM methods, or optimized software implementations. Chord counts alone do not measure runtime. The relevant contribution is a transparent comparison of accuracy mechanisms and their geometric sampling costs.
\begin{thebibliography}{9}
\bibitem{ramanujan} S. Ramanujan, ``Modular equations and approximations to $\pi$,'' \emph{Quarterly Journal of Mathematics} 45 (1914), 350--372.
\bibitem{borwein} J. M. Borwein and P. B. Borwein, \emph{Pi and the AGM}, Wiley, 1987.
\bibitem{ellipse} P. F. Byrd and M. D. Friedman, \emph{Handbook of Elliptic Integrals for Engineers and Physicists}, Springer, 1954.
\bibitem{richardson} L. F. Richardson, ``The approximate arithmetical solution by finite differences of physical problems,'' \emph{Philosophical Transactions of the Royal Society A} 210 (1911), 307--357.
\bibitem{BakerFactorization} C. W. Baker, \emph{Exact Chord Factorization and Universal Asymptotics for Inscribed Ellipse Polygons}, companion manuscript, revised October 2026.
\bibitem{trefethen} L. N. Trefethen and J. A. C. Weideman, ``The exponentially convergent trapezoidal rule,'' \emph{SIAM Review} 56 (2014), 385--458.
\end{thebibliography}
\end{document}
